% GENERATED FILE. Edit canonical lessons, then run npm run generate:books.
% canonical-source-hash: fnv1a64:81e74263456bd20b
% canonical-lessons: IT-C481-fianco, IT-C481-coscia, IT-C481-polpaccio, IT-C481-tallone, IT-C481-alluce

\chapter{Dalla coscia al tallone}
\label{ch:it-dalla-coscia-al-tallone}
\hlchaptermodality{481}

\begin{chapteropening}
\textbf{By the end of this chapter:} \emph{I can name parts of the legs and feet.}

Complete the last lesson of chapter 481: I can name parts of the legs and feet.
\end{chapteropening}

\section[il fianco]{il fianco — hip, side (of the body)}

\label{lesson:IT-C481-fianco}



\begin{quote}
Before the new one: say the Italian for to regret, then the Italian for to let off steam.
\end{quote}

\subsection*{You'll want to know: il fianco}
\textbf{il fianco} — \textquotedblleft{}hip, side (of the body)\textquotedblright{}.

\begin{quote}
Mi fa male il fianco destro.
\end{quote}

\begin{quote}
\emph{My right side hurts.}
\end{quote}

\emph{Di fianco a} — beside: \emph{di fianco alla porta}.

\begin{cousinweb}
\emph{fianco} comes from Frankish \emph{hlanca} \textquotedblleft{}side, hip\textquotedblright{} (EtymDB), a Germanic word, probably through Old French \emph{flanc}.
\end{cousinweb}

\begin{grammarlens}[title={What you've built}]
One more word for telling, explaining and describing at B1.
\end{grammarlens}

\subsection*{Your turn}
\begin{itemize}
\raggedright
  \item \emph{Say it:} \emph{il fianco}
  \item \emph{Say it:} \emph{il fianco}, once more
  \item \emph{Say it:} a sentence of your own with \emph{il fianco}
  \item \emph{Recall:} say the Italian for to regret, then the Italian for to let off steam, then say \emph{il fianco} again
\end{itemize}

\begin{tcolorbox}[breakable,colback=teal!4,colframe=teal!35!black,title={Before you move on}]
What does il fianco mean? (\textquotedblleft{}Hip, side (of the body)\textquotedblright{}.) Say it once more.
\end{tcolorbox}
\section[la coscia]{la coscia — thigh}

\label{lesson:IT-C481-coscia}



\begin{quote}
Before the new one: say the Italian for to let off steam, then the Italian for hip.
\end{quote}

\subsection*{You'll want to know: la coscia}
\textbf{la coscia} — \textquotedblleft{}thigh\textquotedblright{}.

\begin{quote}
Dopo la corsa mi fanno male le cosce.
\end{quote}

\begin{quote}
\emph{After the run my thighs hurt.}
\end{quote}

The plural is \emph{le cosce}. In the kitchen: \emph{una coscia di pollo} — a chicken leg.

\begin{cousinweb}
\emph{coscia} is inherited from Latin \emph{coxa} \textquotedblleft{}hip\textquotedblright{} (Lewis \& Short); in Italian the word moved down to the thigh.
\end{cousinweb}

\begin{grammarlens}[title={What you've built}]
One more word for telling, explaining and describing at B1.
\end{grammarlens}

\subsection*{Your turn}
\begin{itemize}
\raggedright
  \item \emph{Say it:} \emph{la coscia}
  \item \emph{Say it:} \emph{la coscia}, once more
  \item \emph{Say it:} a sentence of your own with \emph{la coscia}
  \item \emph{Recall:} say the Italian for to let off steam, then the Italian for hip, then say \emph{la coscia} again
\end{itemize}

\begin{tcolorbox}[breakable,colback=teal!4,colframe=teal!35!black,title={Before you move on}]
What does la coscia mean? (\textquotedblleft{}Thigh\textquotedblright{}.) Say it once more.
\end{tcolorbox}
\section[il polpaccio]{il polpaccio — calf (of the leg)}

\label{lesson:IT-C481-polpaccio}



\begin{quote}
Before the new one: say the Italian for hip, then the Italian for thigh.
\end{quote}

\subsection*{You'll want to know: il polpaccio}
\textbf{il polpaccio} — \textquotedblleft{}calf (of the leg)\textquotedblright{}.

\begin{quote}
Dopo la gita in montagna ho i polpacci duri.
\end{quote}

\begin{quote}
\emph{After the trip to the mountains my calves are stiff.}
\end{quote}

Built on \emph{polpa} \textquotedblleft{}flesh, pulp\textquotedblright{}.

\begin{cousinweb}
\emph{polpaccio} is built on \emph{polpa} \textquotedblleft{}flesh, pulp\textquotedblright{}, which is inherited from Latin \emph{pulpa} \textquotedblleft{}the fleshy part\textquotedblright{} (Lewis \& Short).
\end{cousinweb}

\begin{grammarlens}[title={What you've built}]
One more word for telling, explaining and describing at B1.
\end{grammarlens}

\subsection*{Your turn}
\begin{itemize}
\raggedright
  \item \emph{Say it:} \emph{il polpaccio}
  \item \emph{Say it:} \emph{il polpaccio}, once more
  \item \emph{Say it:} a sentence of your own with \emph{il polpaccio}
  \item \emph{Recall:} say the Italian for hip, then the Italian for thigh, then say \emph{il polpaccio} again
\end{itemize}

\begin{tcolorbox}[breakable,colback=teal!4,colframe=teal!35!black,title={Before you move on}]
What does il polpaccio mean? (\textquotedblleft{}Calf (of the leg)\textquotedblright{}.) Say it once more.
\end{tcolorbox}
\section[il tallone]{il tallone — heel (of the foot)}

\label{lesson:IT-C481-tallone}



\begin{quote}
Before the new one: say the Italian for thigh, then the Italian for calf.
\end{quote}

\subsection*{You'll want to know: il tallone}
\textbf{il tallone} — \textquotedblleft{}heel (of the foot)\textquotedblright{}.

\begin{quote}
Queste scarpe nuove mi fanno male al tallone.
\end{quote}

\begin{quote}
\emph{These new shoes hurt my heel.}
\end{quote}

The heel of a shoe is \emph{il tacco}.

\begin{cousinweb}
\emph{tallone} is usually traced to Vulgar Latin \emph{tālōnem}, from Latin \emph{tālus} \textquotedblleft{}ankle, heel bone\textquotedblright{} (Lewis \& Short).
\end{cousinweb}

\begin{grammarlens}[title={What you've built}]
One more word for telling, explaining and describing at B1.
\end{grammarlens}

\subsection*{Your turn}
\begin{itemize}
\raggedright
  \item \emph{Say it:} \emph{il tallone}
  \item \emph{Say it:} \emph{il tallone}, once more
  \item \emph{Say it:} a sentence of your own with \emph{il tallone}
  \item \emph{Recall:} say the Italian for thigh, then the Italian for calf, then say \emph{il tallone} again
\end{itemize}

\begin{tcolorbox}[breakable,colback=teal!4,colframe=teal!35!black,title={Before you move on}]
What does il tallone mean? (\textquotedblleft{}Heel (of the foot)\textquotedblright{}.) Say it once more.
\end{tcolorbox}
\section[l'alluce]{l'alluce — big toe}

\label{lesson:IT-C481-alluce}



\begin{quote}
Before the new one: say the Italian for calf, then the Italian for heel.
\end{quote}

\subsection*{You'll want to know: l'alluce}
\textbf{l'alluce} — \textquotedblleft{}big toe\textquotedblright{}.

\begin{quote}
Ho sbattuto l'alluce contro il letto.
\end{quote}

\begin{quote}
\emph{I stubbed my big toe on the bed.}
\end{quote}

The other toes are \emph{le dita dei piedi}.

\begin{cousinweb}
\emph{alluce} is usually traced to Latin \emph{hallux}, a later form of \emph{allex} \textquotedblleft{}the big toe\textquotedblright{} (Lewis \& Short, which also lists \emph{hallus}).
\end{cousinweb}

\begin{grammarlens}[title={What you've built}]
One more word for telling, explaining and describing at B1.
\end{grammarlens}

\subsection*{Your turn}
\begin{itemize}
\raggedright
  \item \emph{Say it:} \emph{l'alluce}
  \item \emph{Say it:} \emph{l'alluce}, once more
  \item \emph{Say it:} a sentence of your own with \emph{l'alluce}
  \item \emph{Recall:} say the Italian for calf, then the Italian for heel, then say \emph{l'alluce} again
\end{itemize}

\begin{tcolorbox}[breakable,colback=teal!4,colframe=teal!35!black,title={Before you move on}]
What does l'alluce mean? (\textquotedblleft{}Big toe\textquotedblright{}.) Say it once more.
\end{tcolorbox}
